\subsection{单变量的全纯函数演算}

定理 5. --- 设 A 是一个幺 Banach 代数，不必交换。设 a 是 A 的元素，z 是 C 的恒等函数在 SpA(a) 的某邻域中的芽。存在唯一的连续保单位元态射 \(\varphi_{a}\)：从 \(\mathcal{O}(\mathrm{Sp}_{\mathrm{A}}(a))\) 到 A 中，使得 \(\varphi_{a}(z)=a\)。

\(\varphi_{a}\) 的像包含在由 \(a\) 生成的 A 的闭全子代数中。特别地，它包含在 \(a\) 的双交换子中。

我们来证明态射 \(\varphi_{a}\) 的存在性。设 B 是由 \(a\) 生成的 A 的闭全子代数。它是交换的，并且 \(\mathrm{Sp}_{\mathrm{B}}(a) = \mathrm{Sp}_{\mathrm{A}}(a)\)（I, p.~5, 第 5 小节）。B 上的全纯函数演算的映射 \(\Theta_{a}\) 是从 \(\mathcal{O}(\mathrm{Sp}_{\mathrm{B}}(a))\) 到 B 的连续保单位元态射，满足 \(\Theta_{a}(z) = a\)（I, p.~51 的定理 1）。\(\Theta_{a}\) 与 B 到 A 中的典范内射的复合态射是一个连续保单位元态射 \(\varphi_{a}\)：从 \(\mathcal{O}(\mathrm{Sp}_{\mathrm{A}}(a))\) 到 A 中，它把 \(z\) 映为 \(a\)。

我们来证明唯一性。设 \(\varphi_{a}^{\prime}\) 是从 \(\mathcal{O}(\mathrm{Sp}_{\mathrm{A}}(a))\) 到 A 的连续保单位元态射，满足 \(\varphi_{a}^{\prime}(z)=a\)。那么 \(\varphi_{a}\) 与 \(\varphi_{a}^{\prime}\) 在 \(\mathrm{Sp}_{\mathrm{A}}(a)\) 的某邻域中的多项式的芽的集合上相一致，因而在 \(\mathrm{Sp}_{\mathrm{A}}(a)\) 的某邻域中的全纯有理分式的芽的集合上相一致。而这些芽在 \(\mathcal{O}(\mathrm{Sp}_{\mathrm{A}}(a))\) 中稠密（I, p.~69, 定理 3）。由此推出 \(\varphi_{a}=\varphi_{a}^{\prime}\)。

\(\varphi_{a}\) 的构造表明其像包含在交换子代数 B 中，而 B 包含在 a 的双交换子中（I, p.~6）。

如果代数 A 的根为零，态射 \(\varphi_{a}\) 的唯一性无需要求连续性即成立（参见 I, p.~40 的命题 9）。一般情形并非如此，参见 G. R. ALLAN, Embedding the algebra of formal power series in a Banach algebra, Proc. London Math. Soc. (3) 25 (1972), 329--340。

对每个 Banach 代数 A、A 的每个元素 \(a\) 以及每个芽 \(f \in \mathcal{O}(\mathrm{Sp}_{\mathrm{A}}(a))\)，我们用 \(f(a)\) 表示定理 5 中的元素 \(\varphi_{a}(f)\)。若 A 是交换 Banach 代数，则元素 \(f(a)\) 与交换 Banach 代数上的全纯函数演算所给出的元素 \(f(a)\) 相一致（I, p.~51 的定理 1）。

设 B 是由 \(a\) 生成的闭全子代数，于是 \(\mathrm{Sp}_{\mathrm{A}}(a) = \mathrm{Sp}_{\mathrm{B}}(a)\)。A 中的元素 \(f(a)\) 属于 B，并且与相对于代数 B 计算的元素 \(f(a)\) 相一致。

命题 7. --- 设 A 与 B 是幺 Banach 代数，\(\varphi\) 是从 A 到 B 的连续保单位元态射。设 \(a \in A\)。则 \(\mathrm{Sp}_{\mathrm{B}}(\varphi(a)) \subset \mathrm{Sp}_{\mathrm{A}}(a)\)，并且有 \(\varphi(f(a)) = f(\varphi(a))\)（对所有 \(f \in \mathcal{O}(\mathrm{Sp}_{\mathrm{A}}(a))\)）。特别地，对每个 \(\chi \in \mathsf{X}(\mathrm{A})\)，有 \(\chi(f(a)) = f(\chi(a))\)。

这由 I, p.~66 的命题 3 得出。

命题 8. --- 设 A 是一个幺 Banach 代数，\(a \in A\)。设 \(f \in \mathcal{O}(\mathrm{Sp}(a))\)。我们有 \(f(\mathrm{Sp}_{\mathrm{A}}(a)) = \mathrm{Sp}_{\mathrm{A}}(f(a))\)。此外，对每个 \(g \in \mathcal{O}(\mathrm{Sp}_{\mathrm{A}}(f(a)))\)，有 \(g \circ f \in \mathcal{O}(\mathrm{Sp}_{\mathrm{A}}(a))\) 且 \(g(f(a)) = (g \circ f)(a)\)。

这由定理 4 得出。

命题 9. --- 设 A 是一个幺 Banach 代数，\(a \in A\)。设 U 是 \(\mathrm{Sp}_{\mathrm{A}}(a)\) 的开邻域，\(f \in \mathcal{O}(\mathrm{U})\)。再设 V 是 \(\mathrm{Sp}_{\mathrm{A}}(a)\) 的包含在 U 中的紧邻域，使得 V 是 U 的一个具有定向边界 \(\partial\mathrm{V}\) 的块。

对每个整数 \(n \geqslant 0\)，映射 \(z \mapsto f(z)(z - a)^{-n-1}\) 在 \(\partial V\) 上连续；微分形式 \(z \mapsto f(z)(z - a)^{-n-1} dz\) 在 \(\partial V\) 上可积，并且有

\[
f ^ {(n)} (a) = \frac {n !}{2 i \pi} \int_ {\partial \mathrm{V}} f (z) (z - a) ^ {- n - 1} d z\tag{7}
\]

其中 \(f^{(n)}\in\mathcal{O}(\mathrm{U})\) 是 \(f\) 的第 n 阶导数。

对 \(n\) 用归纳法。当 \(n = 0\) 时，结论由 I, p.~61 的命题 1 得出。现设命题的断言对整数 \(n \geqslant 0\) 成立。设 \(g \in \mathcal{O}(\mathbf{C} - \mathrm{Sp}_{\mathrm{A}}(a); \mathrm{A})\) 是由 \(g(z) = (z - a)^{-n-1} f(z)\) 定义的全纯函数。微分形式 \(g'(z) dz = dg\) 是 \(\mathrm{C}^1\) 类的；由于块 V 是紧的，Stokes 公式（VAR, R2, p.~47, 11.2.3）蕴涵

\[
\int_ {\partial \mathrm{V}} g ^ {\prime} (z) d z = \int_ {\partial \mathrm{V}} d g = 0.
\]

由于 \(g'(z) = (z - a)^{-n-1} f'(z) - (n+1)(z - a)^{-n-2} f(z)\)，由此推出

\[
\int_ {\partial \mathrm{V}} f ^ {\prime} (z) (z - a) ^ {- n - 1} d z = (n + 1) \int_ {\partial \mathrm{V}} f (z) (z - a) ^ {- n - 2} d z.
\]

把归纳假设应用于 \(f'\)，便得

\[
\frac {2 i \pi}{n !} f ^ {(n + 1)} (a) = (n + 1) \int_ {\partial \mathrm{V}} f (z) (z - a) ^ {- n - 2} d z,
\]

这就是命题关于整数 \(n+1\) 的断言。证明完毕。

命题 10. --- 设 A 是一个幺 Banach 代数，U 是 C 的开子集。

\begin{enumerate}
\def\labelenumi{\alph{enumi})}
\tightlist
\item
  集合 \(\Omega\)——由这样的 \(a \in \mathrm{A}\) 组成：\(\operatorname{Sp}_{\mathrm{A}}(a) \subset \mathrm{U}\)——在 A 中是开的；
\item
  设 \(f \in \mathcal{O}(\mathrm{U})\)。映射 \(a \mapsto f(a)\)（从 \(\Omega\) 到 A 中）是全纯的，特别地是连续的。
\end{enumerate}

设 \(a \in \Omega\)。存在 \(\mathrm{Sp}_{\mathrm{A}}(a)\) 的包含在 U 中的紧邻域 V，它是 U 的一个块（VAR, R2, p.~46 与 p.~47, 11.1.3, d)）。

由于 \(a\) 的预解式在无穷远处趋于 0（I, p.~24 的定理 1, \(c\)）），映射 \(z \mapsto \| (z - a)^{-1}\|\) 在 \(\mathbf{C} - \mathring{\mathrm{V}}\) 上有界。用 M 表示它的上确界。若 \(h \in \mathrm{A}\) 满足 \(\| h\| \leqslant (2\mathrm{M})^{-1}\) 且 \(z \in \mathbf{C} - \mathring{\mathrm{V}}\)，则有

\[
z - (a + h) = (1 - h (z - a) ^ {- 1}) (z - a)
\]

并且 \(\|h(z-a)^{-1}\|\leqslant\frac{1}{2}\)，因而 \(z-(a+h)\) 可逆，其逆满足

\[
(z - (a + h)) ^ {- 1} = (z - a) ^ {- 1} \sum_ {n = 0} ^ {\infty} (h (z - a) ^ {- 1}) ^ {n}\tag{8}
\]

其中 \(\| (h(z - a)^{-1})^n\| \leqslant 2^{-n}\)（I, p.~22 的命题 2）。于是 \(\mathrm{Sp}_{\mathrm{A}}(a + h)\) 包含在 V 中，从而包含在 U 中，这就证明了 \(\Omega\) 在 A 中是开的。

设 \(f \in \mathcal{O}(\mathrm{U})\)。用 \(m\) 表示 \(|f(z)|\) 对 \(z \in \partial \mathrm{V}\) 的上确界。设 \(a \in \mathrm{A}\)。对每个 \(h \in \mathrm{A}\)（满足 \(\| h \| \leqslant (2\mathrm{M})^{-1}\)），有

\[
f (a + h) = \frac {1}{2 i \pi} \int_ {\partial \mathrm{V}} f (z) (z - (a + h)) ^ {- 1} d z
\]

（命题 9）。级数 (8) 在 V 的边界上一致收敛，因而

\[
f (a + h) = \sum_ {n = 0} ^ {+ \infty} f _ {a, n} (h)
\]

其中从 A 到 A 的映射 \(f_{a,n}\) 由下式定义

\[
f _ {a, n} (h) = \frac {1}{2 i \pi} \int_ {\partial \mathrm{V}} f (z) (z - a) ^ {- 1} (h (z - a) ^ {- 1}) ^ {n} d z.
\]

对每个 \(n \in N\)，函数 \(f_{a,n}\) 是 n 次连续齐次多项式函数。此外，由此可得

\[
\left\| f _ {a, n} (h) \right\| \leqslant \frac {m M}{\pi} \left(\int_ {\partial V} \| d z \|\right) 2 ^ {- (n + 1)}
\]

（INT, VI, §2, 第 3 小节, 命题 5）。因此级数 \(\sum_{n} f_{a,n}(h)\) 对 \(\|h\| \leqslant (2M)^{-1}\) 绝对收敛。这就证明了把 \(a\) 映为 \(f(a)\) 的映射在 \(\Omega\) 上全纯（VAR, R1, p.~26, 3.2.1）。

命题 11. --- 设 A 是一个幺 Banach 代数，\(a \in A\)，U 是 \(\mathrm{Sp}_{\mathrm{A}}(a)\) 的开邻域。用 \(\delta\) 表示 \(\mathrm{Sp}_{\mathrm{A}}(a)\) 到 \(\mathbf{C} - \mathbf{U}\) 的距离。设 \(f \in \mathcal{O}(\mathbf{U})\)。

\begin{enumerate}
\def\labelenumi{\alph{enumi})}
\item
  对每个满足 0 \textless{} η \textless{} δ 的实数 η，存在实数 C ≥ 0，使得 \(\|f^{(n)}(a)\| \leqslant \mathrm{Cn}!\eta^{-n}\)（对每个整数 \(n \in N\)）；
\item
  若 \(b \in A\) 与 \(a\) 交换且 \(\varrho(b) < \delta\)，则 \(\mathrm{Sp}_{\mathrm{A}}(a + b) \subset \mathrm{U}\)，并且
\end{enumerate}

\[
f (a + b) = \sum_ {n = 0} ^ {\infty} \frac {f ^ {(n)} (a)}{n !} b ^ {n},
\]

其中级数绝对收敛。

设 \(\eta\) 是满足 \(0 < \eta < \delta\) 的实数。记 \(\varepsilon = \delta - \eta > 0\)。设 K 是 \(\mathrm{Sp}_{\mathrm{A}}(a)\) 的紧邻域，它由 C 中到 \(\mathrm{Sp}_{\mathrm{A}}(a)\) 的距离 \(\leqslant \varepsilon / 2\) 的点组成。由于 \(f\) 在每个以属于 K 的点为中心、半径为 \(\eta + \varepsilon / 2\) 的开圆盘中全纯，由 Cauchy 不等式（VAR, R1, p.~29, 3.3.4），存在实数 \(C \geqslant 0\)，使得

\[
\sup _ {z \in K} \frac {| f ^ {(n)} (z) |}{n !} \leqslant \frac {C}{\eta^ {n}}
\]

对每个整数 \(n \geqslant 0\) 成立。于是把 I, p.~61 的命题 1 应用于 \(f^{(n)}\) 以及包含在 K 中的一个块 V，便得断言 a)。

设 \(b\) 是 A 中与 \(a\) 交换的元素，满足 \(\varrho(b) < \delta\)。用由 \(a\) 与 \(b\) 生成的闭全子代数 B 代替 A，该子代数满足 \(\mathrm{Sp}_{\mathrm{A}}(a) = \mathrm{Sp}_{\mathrm{B}}(a)\) 与 \(\mathrm{Sp}_{\mathrm{A}}(a + b) = \mathrm{Sp}_{\mathrm{B}}(a + b)\)，为证明 \(b\))，就可化为 A 为交换的情形。

由于 \(\varrho(b) < \delta\)，可以选取 \(\eta\) 使得 \(\varrho(b) < \eta < \delta\)。设 \(V_1\) 是 \(\mathbf{C}\) 中到 \(\mathrm{Sp}_{\mathrm{A}}(a)\) 的距离 \(< \delta - \eta\) 的点组成的集合，\(V_2\) 是以 0 为中心、半径为 \(\eta\) 的开圆盘（在 \(\mathbf{C}\) 中）。设 \(g\) 是映射 \((z_1, z_2) \mapsto z_1 + z_2\)：从 \(V_1 \times V_2\) 到 U 中。那么 \(h = f \circ g\) 是映射 \((z_1, z_2) \mapsto f(z_1 + z_2)\)：从 \(V_1 \times V_2\) 到 \(\mathbf{C}\) 中。我们有 \(\mathrm{Sp}_{\mathrm{A}}^{2}(a, b) \subset V_1 \times V_2\)，因而 \(\mathrm{Sp}_{\mathrm{A}}(a + b) \subset U\)（参见 I, p. 67 的推论 2），并且 \(f(a + b) = h(a, b)\)（由 I, p.~72 的定理 4）。现在，在空间 \(\mathcal{O}(V_1 \times V_2)\) 中，有

\[
h (z _ {1}, z _ {2}) = \sum_ {n \geqslant 0} \frac {f ^ {(n)} (z _ {1})}{n !} z _ {2} ^ {n},
\]

因此级数

\[
\sum_ {n \geqslant 0} \frac {f ^ {(n)} (a)}{n !} b ^ {n}
\]

在 A 中收敛，其和为 \(h(a,b)=f(a+b)\)。此外，由断言 (a)，这个级数绝对收敛。

